% ============================== CHAPTER 3 ==============================
\chapterdivider{3}{PROPOSED METHODOLOGY}
\chapter{Proposed Methodology}

\section{Introduction}
The proposed system is an open, measured study of pseudonym resolution in the protocol of~\cite{base}. It consists of a reference implementation of the protocol whose server can be configured with any of seven resolvers, instrumentation that counts every hash evaluation the server makes, a network model that prices those counts with the paper's own constants, and symbolic models that check the protocol's security properties for every correct resolver.

The methodology follows the order in which each result depends on the previous one:
\begin{enumerate}
\item \textbf{Specification and conformance.} Implement Section~III of~\cite{base} and fix its message layout with independent reference vectors, in Python and in C.
\item \textbf{Gap analysis.} Using only the specification and the vectors, show that resolution is required and identify the choices it leaves open.
\item \textbf{Design space.} Implement seven resolution strategies along those choices, each with a closed-form cost.
\item \textbf{Measurement and validation.} Count probes for honest and forged pseudonyms for $N=1$ to $10{,}000$ and check the counts against the closed forms.
\item \textbf{Network evaluation.} Feed the measured cost model into an OMNeT++ model of the paper's Table~II network and rank the strategies.
\item \textbf{Failure modes and formal verification.} Test desynchronisation and replay, and verify secrecy, agreement and unlinkability in ProVerif.
\end{enumerate}
Each phase serves one of the objectives of Section~1.3, and every number reported later can be regenerated from the repository.

\section{Requirement Gathering}
The project is a research study of a published protocol, so its requirements were gathered by \emph{document analysis} rather than by surveys or interviews. The team studied the base paper~\cite{base} equation by equation (Eqs.~8--24, Tables~II--VI) and recorded, for every equation, the code that implements it and the evidence that checks it (the equation-to-code matrix in \texttt{docs/conformance.md}). Every decision that the paper does not fix (epoch semantics, search window, failure threshold, nonce length, persistence of the epoch) was written into an \emph{assumption register} together with its basis. The RFID literature reviewed in Chapter~2 supplied the candidate resolution strategies and their known costs. Requirements were refined through review meetings with the project guide.

\subsection{Functional Requirements}
\begin{enumerate}[label=FR-\arabic*]
\item \textbf{Enrolment.} A trusted authority shall derive each entity's long-term secret $\lambda_X=H_1(ID_X\|s)$ (Eq.~8), issue credentials to devices and servers, and share peer secrets so that each pair can derive $k^\star$ (Eq.~9).
\item \textbf{Request.} A device shall build $M_1=\langle PID,\mathrm{AEAD}_{k^\star}(PID\|N_i\|t)\rangle$ with $PID=H_2(\lambda_V\|t)$ (Eqs.~11--15).
\item \textbf{Resolution.} A server shall map an incoming $PID$ to an enrolled device and an epoch using a configurable resolver, and shall count the hash evaluations it makes.
\item \textbf{Verification and response.} The server shall decrypt and verify $M_1$, check that the decrypted $PID$ and $t$ match, and return $M_2=\langle ID_R,\mathrm{AEAD}_{k^\star}(N_j\|t;\,PID\|N_i)\rangle$ (Eqs.~17--19).
\item \textbf{Key agreement.} Both sides shall derive the same session key $K=\mathrm{KDF}(k^\star\|N_i\|N_j\|t)$ (Eq.~16).
\item \textbf{Replay and abuse handling.} The server shall reject a replayed $M_1$, and both sides shall block a peer after three failed attempts.
\item \textbf{Strategies.} The server shall support strategies O1--O4 and TM; strategies CR and MW shall be modelled at the resolution level.
\item \textbf{Measurement.} Scripts shall measure probes per honest and forged message and the resolution state for $N=1$ to $10{,}000$, and write a cost model for the simulator.
\item \textbf{Simulation and verification.} An OMNeT++ model shall price the cost model under the paper's network parameters, and ProVerif models shall check the security properties.
\item \textbf{Demonstration.} Command-line tools shall enrol entities, run a server, authenticate a device and demonstrate replay rejection; the C device port shall interoperate with the Python server.
\end{enumerate}

\subsection{Non Functional Requirements}
\begin{itemize}
\item \textbf{Security.} Standard primitives only: AES-256-GCM~\cite{gcm}, HKDF-SHA256~\cite{hkdf}, SHA-256~\cite{fips180}; nonces from a cryptographically secure generator; constant-time comparison of each candidate pseudonym; secrets never printed (only a fingerprint of the session key).
\item \textbf{Conformance.} The Python and C implementations shall reproduce reference vectors generated by code that shares nothing with them, byte for byte.
\item \textbf{Reliability and quality.} At least 90\% statement coverage of the protocol and crypto packages, enforced in continuous integration on Python~3.11 and~3.12.
\item \textbf{Performance.} Resolution cost shall be reported in probes and priced at the paper's $T_H=0.32$\,ms, so that results are comparable with the paper's 23.68\,ms per session.
\item \textbf{Reproducibility.} Pinned tool versions (OMNeT++~6.4.0, ProVerif~2.05), committed data and logs, and one command per result.
\item \textbf{Portability.} The device side shall run on an 80\,MHz ESP8266 without hardware cryptography.
\item \textbf{Scalability.} Measurements shall cover up to $10{,}000$ devices per server and simulations up to $1{,}000$ devices per network.
\end{itemize}

\subsection{User Requirements}
\begin{itemize}
\item \textbf{Protocol designers and researchers} need to know which properties of the protocol depend on the resolution choice, and what term the cost model is missing.
\item \textbf{Implementers and deployers} need a recommendation: which strategy to choose for their network size, memory budget and threat model, and what each choice costs.
\item \textbf{Reviewers and examiners} need to reproduce every number from the released code and data, and to inspect results without installing the simulator or verifier.
\item \textbf{Device owners} need devices that do only symmetric work online and that are not locked out by an attacker who drops a few messages.
\end{itemize}

\section{Proposed Algorithms}
\subsection*{The base protocol}
In the offline phase the trusted authority computes, for each entity $X$ and a master secret $s$,
\begin{equation}
\lambda_X = H_1(ID_X\|s), \qquad k^\star = \mathrm{KDF}(\lambda_V\|\lambda_R),
\end{equation}
so that device $V$ and server $R$ share the pairwise key $k^\star$. Online, in epoch $t$, the device samples a nonce $N_i$ and sends
\begin{equation}
M_1=\langle PID,\;\mathrm{AEAD}_{k^\star}(PID\|N_i\|t)\rangle,\qquad PID=H_2(\lambda_V\|t).
\end{equation}
The server resolves $PID$ to $(V,t)$, decrypts and verifies $M_1$, samples $N_j$, and returns
\begin{equation}
M_2=\langle ID_R,\;\mathrm{AEAD}_{k^\star}(N_j\|t;\;PID\|N_i)\rangle ,
\end{equation}
where $PID\|N_i$ is associated data. Both sides derive
\begin{equation}
K=\mathrm{KDF}(k^\star\|N_i\|N_j\|t).
\end{equation}
We instantiate AEAD as AES-256-GCM, KDF as HKDF-SHA256 and $H$ as SHA-256. Table~\ref{tab:notation} summarises the notation.

\begin{table}[H]
\centering
\caption{Notation}
\label{tab:notation}
\small
\begin{tabular}{|l|p{0.72\textwidth}|}
\hline
\textbf{Symbol} & \textbf{Meaning}\\ \hline
$V$, $R$ & Device and server\\ \hline
$\lambda_X$ & Long-term secret of entity $X$ (Eq.~8)\\ \hline
$k^\star$ & Pairwise key of a device--server pair (Eq.~9)\\ \hline
$t$, $PID$ & Epoch, and pseudonym $H_2(\lambda_V\|t)$ (Eq.~11)\\ \hline
$N_i$, $N_j$ & Device and server nonces (128 bits)\\ \hline
$N$, $n$ & Devices in the network; devices enrolled at one server\\ \hline
$L$, $F$, $W$ & Backward and forward epoch tolerance ($L=16$, $F=8$) and window $W=L+F+1=25$\\ \hline
$W_t$ & Clock-epoch drift window ($W_t=3$)\\ \hline
$m$ & Device lifetime in epochs for TM ($m=128$); $S=Nm$\\ \hline
$\delta$, $d$ & Branching factor and depth of the key tree ($\delta=32$)\\ \hline
$T_H$ & Cost of one hash, 0.32\,ms (Table~III of~\cite{base}); one \emph{probe} is one $H_2$ evaluation\\ \hline
\end{tabular}
\end{table}

\subsection*{Why resolution is required}
Each step of the argument uses only the specification or a fact confirmed by the reference vectors:
\begin{enumerate}
\item $M_1$ carries $PID$ in clear; $t$ and $N_i$ are only inside the ciphertext under $k^\star$ (Eqs.~13--15, confirmed byte for byte).
\item $k^\star=\mathrm{KDF}(\lambda_V\|\lambda_R)$ is specific to the device (Eqs.~9, 12).
\item So the server cannot decrypt $M_1$ until it knows which device sent it.
\item The paper requires that $PID$ cannot be linked to $V$ without $\lambda_V$ (its Section~III-J), so the identity cannot be read from $M_1$ either.
\item So every correct server must map $PID$ to $(V,t)$ by computing pseudonyms from the secrets it holds.
\end{enumerate}
Two choices are left open and together decide the cost: the \emph{epoch semantics} (a per-device counter or a clock) and the \emph{resolution method} (an on-demand search or a table prepared in advance). Figure~\ref{fig:flow} shows where the step sits.

\begin{figure}[H]
\centering
\begin{tikzpicture}[font=\small, >=Stealth]
\node[draw, rounded corners, minimum width=2.2cm, minimum height=7mm, fill=gray!8] (D) {Device $V$};
\node[draw, rounded corners, minimum width=2.2cm, minimum height=7mm, fill=gray!8, right=8.2cm of D] (S) {Server $R$};
\draw[gray, dashed] (D.south) -- ++(0,-6.0);
\draw[gray, dashed] (S.south) -- ++(0,-6.0);
\draw[->, thick] ($(D.south)+(0,-0.6)$) -- node[above]{$M_1=\langle PID,\;\mathrm{AEAD}_{k^\star}(PID\|N_i\|t)\rangle$} ($(S.south)+(0,-0.6)$);
\node[draw=orange!80!black, fill=orange!10, align=left, anchor=north east, text width=6.0cm] (Rz) at ($(S.south)+(-0.15,-1.0)$)
 {\textbf{Resolution (unspecified)}\\ find $(V,t)$ with $H_2(\lambda_V\|t)=PID$ over $N$ devices and $W$ epochs: scan, table, tree or trade-off};
\node[draw, fill=gray!8, align=left, anchor=north east, text width=6.0cm] (Vf) at ($(Rz.south east)+(0,-0.2)$)
 {decrypt with $k^\star_V$, verify $\mu_1$\\ (Table~V of~\cite{base}: $2T_H+2T_{SE/D}$)};
\draw[<-, thick] ($(D.south)+(0,-5.6)$) -- node[above]{$M_2=\langle ID_R,\;\mathrm{AEAD}_{k^\star}(N_j\|t;\,PID\|N_i)\rangle$} ($(S.south)+(0,-5.6)$);
\end{tikzpicture}
\caption{Where pseudonym resolution sits in the exchange. The server needs $k^\star_V$ to open $M_1$, but $V$ is hidden behind $PID$.}
\label{fig:flow}
\end{figure}

\subsection*{The seven resolution strategies}
The cost of resolution is counted in \emph{probes}, one probe being one evaluation of $H_2$ (or the equivalent keyed hash). The first five strategies keep the wire format of $M_1$; CR and MW need a fresh value $r$ in clear and therefore change $M_1$.

\textbf{O1 -- counter epochs, linear scan} (the behaviour of our original implementation; the OSK search~\cite{ohkubo}). For each enrolled device the server probes the epochs from $\ell-L$ to $\ell+F$ around the last accepted epoch $\ell$.
\begin{lstlisting}[style=pseudo]
function ResolveO1(PID):
  for each enrolled device V with secret $\lambda_V$ and last accepted epoch $\ell_V$ do
    for t = $\ell_V - L$ to $\ell_V + F$ do
      if $H_2(\lambda_V \| t)$ = PID then            // constant-time compare
        if t $\le \ell_V$ then raise ReplayDetected
        return (V, t)
  raise UnknownPeer                          // a forged PID costs N$\cdot$W probes
\end{lstlisting}

\textbf{O2 -- counter epochs, precomputed table.} The server keeps a dictionary from each device's next $F$ pseudonyms to $(V,t)$. Resolution is a lookup; after an epoch is accepted the table slides, which costs one probe in the normal case.
\begin{lstlisting}[style=pseudo]
setup: for each V: for t = $\ell_V$+1 to $\ell_V$+F: T[$H_2(\lambda_V\|t)$] = (V, t)
function ResolveO2(PID):
  if PID in T then return T[PID] else raise UnknownPeer   // no hash on a miss
function AcceptO2(V, t):
  remove entries of V with epoch $\le$ t;  add $H_2(\lambda_V\|t')$ for new t' up to t+F
\end{lstlisting}

\textbf{O3 -- clock epochs, linear scan.} The epoch is $t=\lfloor\text{clock}/\text{epoch length}\rfloor$; the server probes every device for each epoch in a drift window of $W_t=3$ epochs, and keeps a cache of accepted nonces to detect replays.

\textbf{O4 -- clock epochs, per-epoch table} (the time-based idea of YA-TRAP~\cite{tsudik}). At each epoch boundary the server adds every device's pseudonym for the new epoch ($N$ probes per epoch) and drops expired ones; resolution is a lookup with no probes.

\textbf{TM -- rainbow time--memory trade-off}~\cite{avoine05,ado05,hellman,oechslin}. Over the space of $S=Nm$ points (device, epoch) with lifetime $m=128$, the server builds six rainbow tables of $2S^{2/3}$ chains of length $\lceil S^{1/3}\rfloor$ and stores only chain end points. A lookup regenerates chains from the received $PID$.

\textbf{CR -- randomized hash lock}~\cite{weis}. The device sends $(r, H(k\|r))$ for a fresh 128-bit $r$; the server tries every key.

\textbf{MW -- key tree}~\cite{molnar}. Keys form a tree with branching factor $\delta=32$; the device sends one keyed hash per level and the server descends level by level, trying at most $\delta$ keys per level.

Table~\ref{tab:design} gives the closed-form cost of each strategy, which the measurements in Chapter~4 are checked against.

\begin{table}[H]
\centering
\caption{Design space of resolution strategies. One table entry is charged at 40\,B; ``probes'' are expected $H_2$ evaluations per message in steady state.}
\label{tab:design}
\scriptsize
\setlength{\tabcolsep}{3pt}
\renewcommand{\arraystretch}{1.4}
\newcolumntype{Y}{>{\raggedright\arraybackslash}X}
\begin{tabularx}{\textwidth}{|l|>{\raggedright\arraybackslash}p{1.45cm}|Y|Y|>{\raggedright\arraybackslash}p{1.25cm}|Y|>{\raggedright\arraybackslash}p{1.3cm}|>{\raggedright\arraybackslash}p{1.15cm}|}
\hline
 & \textbf{Epoch / $M_1$} & \textbf{Probes, honest} & \textbf{Probes, forged} & \textbf{Pre\-compute} & \textbf{Resolution state} & \textbf{Dropped $M_1$} & \textbf{Link\-ability}\\ \hline
O1 & counter / unchanged & $\frac{(N-1)W}{2}+L+2$ & $NW$ & none & $8N$\,B & lockout after $F$ & none\\ \hline
O2 & counter / unchanged & 1 (slide) & 0 & $NF$ once & $(40F+8)N$\,B & lockout after $F$ & none\\ \hline
O3 & clock / unchanged & $\frac{(N-1)W_t}{2}+1$ & $NW_t$ & none & nonce cache & none & same epoch\\ \hline
O4 & clock / unchanged & 0 & 0 & $N$ per epoch & $40W_tN$\,B + nonce cache & none & same epoch\\ \hline
TM & counter / unchanged & $\approx0.5$--$0.7\,S^{2/3}$ (measured) & $\approx4\,S^{2/3}$ (measured) & $\approx12S$ once & $\approx50\,S^{2/3}$\,B & none until $t=m$ & none\\ \hline
CR & none / $+r$ (128\,b) & $\frac{N+1}{2}$ & $N$ & none & none & none & none\\ \hline
MW & none / $+r$, $+(d-1)$ tags & $\le\delta d$ & $\le\delta$ & none & $32\times$ internal nodes & none & after compromise\\ \hline
\end{tabularx}
\end{table}

\subsection*{Evaluation method}
The network evaluation follows the paper's conditions and changes only the resolution term:
\begin{itemize}
\item \textbf{Simulator and environment:} OMNeT++~6.4~\cite{omnet} (Cmdenv) with Table~II of~\cite{base}: a $1500\times1500$\,m area, 10\,Mbps, 500\,m range, 75.6\,mW, a 15\,$\mu$s quantum, and servers at 5\% of the devices.
\item \textbf{Protocol cost:} Tables~III--V of~\cite{base}: $2T_H+2T_{SE/D}=11.84$\,ms per side, 1,024 bits over four messages, and $T_H=0.32$\,ms per probe.
\item \textbf{Resolution cost:} for each server, the probes for its number of enrolled devices $n$ are read from the measured cost model (exact for O1--O4 and CR, interpolated between measured points for TM and MW). CR and MW also add their extra $M_1$ bits.
\item \textbf{Assumptions} (not fixed by the paper): 5\% adversaries, each sending forged $M_1$ at 1\,Hz; a 10\,s authentication window; each server enrols every device in its range; first-come-first-served processing; one O4 table extension per server per window.
\item \textbf{Runs:} five seeds per point, $N\in\{1,10,100,200,\dots,1000\}$, with and without adversaries: 960 runs in total, reported as means.
\end{itemize}
For formal verification, the exchange is modelled in ProVerif~2.05~\cite{proverif} in the symbolic model with an active attacker, with unbounded servers, honest and corrupt devices and sessions. The server may resolve a pseudonym to \emph{any} enrolled device and must then decrypt and check $PID=H_2(\lambda_V\|t)$, so every property that holds, holds for every resolver that returns a correct device. Unlinkability is checked as observational equivalence~\cite{blanchet08} in the style of~\cite{arapinis}.

\section{Gantt chart for Timeline}\label{sec:gantt}
Figure~\ref{fig:gantt} reconstructs the project schedule from the team's account and the confirmed academic milestones. The base paper became available on 6 February 2026. The Semester~VI work covered literature review, understanding the specification and environment setup, particularly between Reviews~6.1 and~6.2. The implementation and evaluation were completed during the interval between Reviews~7.1 and~7.2. Activity bars show approximate month-level periods; the internal allocation of implementation, simulation, verification and writing is retrospective, rather than a record of exact start and finish dates. The dated milestones are confirmed by the team. Repository commit dates are not used as activity dates.

\begin{figure}[H]
\centering
\begin{ganttchart}[
  x unit=0.81cm, y unit chart=0.48cm, y unit title=0.48cm,
  hgrid, vgrid,
  bar height=0.48, bar/.append style={fill=blue!35, draw=blue!60!black},
  milestone/.append style={fill=orange!80!black, draw=orange!50!black},
  bar label font=\scriptsize, milestone label font=\scriptsize,
  bar label node/.append style={text width=5.1cm, align=right},
  milestone label node/.append style={text width=5.1cm, align=right},
  title label font=\scriptsize\bfseries, title height=1,
  title list options={var=\x}
]{1}{10}
\gantttitle{2026 --- approximate activity periods}{10} \\
\gantttitlelist{Feb,Mar,Apr,May,Jun,Jul,Aug,Sep,Oct,Nov}{1} \\
\ganttmilestone{Base paper available: 6 Feb}{1} \\
\ganttbar{Initial study and synopsis preparation}{1}{2} \\
\ganttmilestone{Synopsis submitted: 5 Mar}{2} \\
\ganttmilestone{Review 6.1: 17 Mar}{2} \\
\ganttbar{Literature, protocol study and setup}{2}{3} \\
\ganttmilestone{Review 6.2: 15 Apr}{3} \\
\ganttmilestone{Logbook and attendance report: 28 Apr}{3} \\
\ganttmilestone{Review 7.1: 21 Aug}{7} \\
\ganttbar{Python protocol core and demo}{7}{8} \\
\ganttbar{C port and tests on two ESP8266 boards}{8}{9} \\
\ganttbar{Resolvers and probe measurements}{8}{9} \\
\ganttbar{OMNeT++ network evaluation}{9}{9} \\
\ganttbar{ProVerif models and analysis}{9}{9} \\
\ganttbar{Conformance tests, results and paper}{8}{9} \\
\ganttmilestone{Review 7.2: 27 Oct}{9} \\
\ganttbar{Report preparation and final documentation}{9}{10}
\end{ganttchart}
\caption{Retrospective Gantt chart: approximate activity periods and confirmed milestones. Blank months carry no inferred activity dates.}
\label{fig:gantt}
\end{figure}

\section{Architectural Diagram}
Figure~\ref{fig:arch} shows the layered architecture of the system. The \emph{crypto} layer wraps the primitives and contains no protocol logic. The \emph{protocol} layer contains transport-agnostic state machines for the trusted authority, device and server, the message formats and the pluggable resolvers. The \emph{storage} layer holds the session log (replay window and block-list) and a demonstration key store. The \emph{demo} layer adds length-prefixed TCP framing and the command-line tools. The C device core reuses the same wire format. The \emph{evaluation} layer measures, simulates and verifies, and feeds the figures and tables of the paper and of this report.

\begin{figure}[H]
\centering
\resizebox{\textwidth}{!}{%
\begin{tikzpicture}[font=\small, >=Stealth,
  layer/.style={draw, rounded corners=2pt, minimum height=9mm, text width=12.6cm, align=center, fill=#1},
  layer/.default=gray!6]
\node[layer=green!8] (eval) {\textbf{Evaluation}\\ \texttt{resolution\_costs.py}, \texttt{rfid\_schemes.py}, \texttt{desync\_check.py} $\cdot$ OMNeT++ \texttt{ResolveNet} $\cdot$ ProVerif models $\cdot$ \texttt{make\_figs.py}};
\node[layer=yellow!10, below=3mm of eval] (demo) {\textbf{Demo / transport}\\ \texttt{hea-ta}, \texttt{hea-server}, \texttt{hea-device} $\cdot$ length-prefixed TCP framing (\texttt{\_transport.py})};
\node[layer=blue!6, below=3mm of demo] (proto) {\textbf{Protocol}\\ \texttt{ta.py} (TrustedAuthority) $\cdot$ \texttt{entity.py} $\cdot$ \texttt{device.py} $\cdot$ \texttt{server.py} $\cdot$ \texttt{resolution.py} (O1--O4, TM) $\cdot$ \texttt{messages.py}};
\node[layer=orange!8, below=3mm of proto] (stor) {\textbf{Storage}\\ \texttt{session\_log.py} (epoch window, replay, 3-attempt block-list) $\cdot$ \texttt{keystore.py} (demo-grade)};
\node[layer=red!6, below=3mm of stor] (cry) {\textbf{Crypto}\\ \texttt{hashes.py} ($H_1,H_2,H_3$) $\cdot$ \texttt{kdf.py} (HKDF-SHA256) $\cdot$ \texttt{aead.py} (AES-256-GCM) $\cdot$ \texttt{nonce.py} $\cdot$ \texttt{curve.py} (P-256, offline only)};
\node[draw, rounded corners=2pt, fill=purple!6, text width=3.0cm, align=center, minimum height=2.6cm, right=4mm of proto.east, yshift=-6mm] (c) {\textbf{C device core}\\ \texttt{hea\_proto.c}\\ BearSSL\\ ESP8266 sketch\\ host build};
\draw[<->, thick] (c.west) -- ++(-0.4,0);
\end{tikzpicture}}
\caption{Layered architecture of the implementation and evaluation}
\label{fig:arch}
\end{figure}

\subsection{UML Diagrams}
\textbf{Use-case diagram.} Figure~\ref{fig:usecase} shows the actors and the services of the system. The trusted authority enrols entities offline. A device requests authentication and derives the session key; the server resolves the pseudonym, verifies the request and responds. An adversary may replay, forge or drop messages. The researcher runs the measurements, simulation and verification.

\newcommand{\actor}[2]{%
  \begin{scope}[shift={(#1)}]
  \draw (0,0.55) circle (0.17);
  \draw (0,0.38) -- (0,-0.05) (-0.25,0.25) -- (0.25,0.25) (0,-0.05) -- (-0.2,-0.4) (0,-0.05) -- (0.2,-0.4);
  \node[below, align=center, font=\footnotesize] at (0,-0.42) {#2};
  \end{scope}}
\begin{figure}[H]
\centering
\begin{tikzpicture}[font=\footnotesize, uc/.style={draw, ellipse, minimum width=3.6cm, minimum height=0.75cm, align=center, fill=blue!4}]
\draw[rounded corners] (-3.0,-4.9) rectangle (3.0,3.4);
\node[font=\small\bfseries] at (0,3.1) {Authentication system};
\node[uc] (u1) at (0,2.45) {Enrol entities (Eqs.~8--9)};
\node[uc] (u2) at (0,1.5) {Send authentication request $M_1$};
\node[uc] (u3) at (0,0.55) {Resolve pseudonym};
\node[uc] (u4) at (0,-0.4) {Verify $M_1$, send $M_2$};
\node[uc] (u5) at (0,-1.35) {Derive session key $K$};
\node[uc] (u6) at (0,-2.3) {Reject replay / block peer};
\node[uc] (u7) at (0,-3.25) {Measure probes, simulate};
\node[uc] (u8) at (0,-4.2) {Verify properties (ProVerif)};
\actor{-5.2,2.2}{Trusted\\Authority}
\actor{-5.2,-0.2}{Device}
\actor{-5.2,-2.9}{Adversary}
\actor{5.2,0.3}{Server}
\actor{5.2,-3.5}{Researcher}
\draw (-4.95,2.4) -- (u1.west);
\draw (-4.95,0.0) -- (u2.west); \draw (-4.95,0.0) -- (u5.west);
\draw[dashed] (-4.95,-2.7) -- (u2.west); \draw[dashed] (-4.95,-2.7) -- (u6.west);
\draw (4.95,0.5) -- (u3.east); \draw (4.95,0.5) -- (u4.east); \draw (4.95,0.5) -- (u5.east); \draw (4.95,0.5) -- (u6.east); \draw (4.95,0.5) -- (u1.east);
\draw (4.95,-3.3) -- (u7.east); \draw (4.95,-3.3) -- (u8.east);
\end{tikzpicture}
\caption{Use-case diagram}
\label{fig:usecase}
\end{figure}

\textbf{Class diagram.} Figure~\ref{fig:class} shows the main classes of the Python implementation. \texttt{Device} and \texttt{Server} share the \texttt{Entity} base class, which caches pairwise keys and implements the coverage and access-control gates (Eqs.~10, 23). A \texttt{Server} owns a \texttt{SessionLog} and one \texttt{Resolver}; the resolver subclasses implement the strategies O1--O4 and TM.

\begin{figure}[H]
\centering
\begin{tikzpicture}[font=\scriptsize, >=Stealth,
  uml/.style={rectangle split, rectangle split parts=3, draw, fill=yellow!6, text width=#1, align=left, inner sep=2.5pt, anchor=north},
  uml/.default=3.3cm,
  umls/.style={draw, fill=yellow!6, text width=#1, align=center, inner sep=2.5pt, anchor=north},
  umls/.default=2.2cm,
  inh/.style={-{Triangle[open,length=2.5mm,width=2.5mm]}},
  comp/.style={{Diamond[length=2.6mm,width=2mm]}-}]
\node[uml] (ta) at (-5.0,0) {\centering\textbf{TrustedAuthority}\nodepart{second}master secret $s$\nodepart{third}enroll(identity, role)\\ enroll\_batch(manifest)\\ share\_peer\_keys(cred, peers)};
\node[uml] (en) at (0,0) {\centering\textbf{Entity}\nodepart{second}credential, session\_log\nodepart{third}get\_pairwise\_root(peer)\\ coverage\_ok(peer)\\ access\_control\_gate(peer, ok)};
\node[uml] (sl) at (5.0,0) {\centering\textbf{SessionLog}\nodepart{second}last accepted epoch, failures\nodepart{third}epoch\_search\_window(peer)\\ is\_stale\_epoch(peer, t)\\ accept\_epoch(peer, t)\\ record\_failure(peer)\\ is\_blocked(peer)};
\node[uml=2.9cm] (cr) at (-5.6,-3.4) {\centering\textbf{Credential}\nodepart{second}identity, role, lam\\ peer\_secrets\nodepart{third}to\_dict(), from\_dict()};
\node[uml] (dv) at (-1.4,-3.4) {\centering\textbf{Device}\nodepart{second}epoch per peer\nodepart{third}current\_pseudonym(peer)\\ build\_auth\_request(peer)\\ complete\_auth(M2)};
\node[uml] (sv) at (2.6,-3.4) {\centering\textbf{Server}\nodepart{second}resolver\nodepart{third}resolve\_pseudonym(PID)\\ handle\_auth\_request(M1)};
\node[uml=3.0cm] (msg) at (-5.5,-6.4) {\centering\textbf{Message1 / Message2}\nodepart{second}PID or server\_id, nonce, ct, tag\nodepart{third}to\_bytes(), from\_bytes()};
\node[uml=3.6cm] (rs) at (2.6,-6.4) {\centering\textit{\textbf{Resolver}}\nodepart{second}probes, precompute\_probes\nodepart{third}attach(secrets, log)\\ resolve(PID) $\to$ (peer, t)\\ accept(peer, t, $N_i$)\\ storage\_bytes()};
\node[umls] (o1) at (-3.9,-10.0) {\textbf{LinearScan}\\ (O1)};
\node[umls] (o2) at (-1.4,-10.0) {\textbf{CounterTable}\\ (O2)};
\node[umls=2.5cm] (ck) at (1.4,-10.0) {\textbf{\_ClockResolver}\\ nonce cache};
\node[umls=2.6cm] (tm) at (4.6,-10.0) {\textbf{TimeMemoryTable}\\ (TM)};
\node[umls] (o3) at (0.1,-11.6) {\textbf{ClockScan}\\ (O3)};
\node[umls] (o4) at (2.7,-11.6) {\textbf{ClockTable}\\ (O4)};
\draw[inh] (dv.north) -- ++(0,0.5) -| ([xshift=-6mm]en.south);
\draw[inh] (sv.north) -- ++(0,0.5) -| ([xshift=6mm]en.south);
\draw[comp] (en.east) -- node[above]{1} (sl.west);
\draw[comp] (sv.south) -- node[right]{1} (rs.north);
\draw[->, dashed] (ta.south) -- node[left]{issues} (cr.north);
\draw[->, dashed] (dv.west) -- node[above]{uses} (cr.east |- dv.west);
\draw[->, dashed] (dv.south) |- node[pos=0.75, below]{builds / parses} (msg.east);
\coordinate (bus) at (0,-9.4);
\draw (o1.north) |- (bus -| o1.north) -- (bus -| tm.north);
\draw (o2.north) -- (o2.north |- bus);
\draw (ck.north) -- (ck.north |- bus);
\draw (tm.north) -- (tm.north |- bus);
\draw[inh] (rs.south |- bus) -- (rs.south);
\coordinate (bus2) at (0,-11.25);
\draw (o3.north) -- (o3.north |- bus2) -- (o4.north |- bus2) -- (o4.north);
\draw[inh] (ck.south |- bus2) -- (ck.south);
\end{tikzpicture}
\caption{Class diagram of the Python implementation}
\label{fig:class}
\end{figure}

\textbf{Sequence diagram.} Figure~\ref{fig:seq} shows one successful authentication, including the resolver and the session log inside the server.

\begin{figure}[H]
\centering
\begin{tikzpicture}[font=\scriptsize, >=Stealth]
\foreach \x/\n/\l in {0/dev/Device, 4.2/srv/Server, 7.6/res/Resolver, 10.6/log/SessionLog}{
  \node[draw, fill=gray!8, minimum width=2.0cm, minimum height=6mm] (\n) at (\x,0) {\l};
  \draw[dashed, gray] (\n.south) -- ++(0,-9.4);}
\def\msg#1#2#3#4{\draw[->] (#1,#3) -- node[above, align=center]{#4} (#2,#3);}
\def\ret#1#2#3#4{\draw[->, dashed] (#1,#3) -- node[above, align=center]{#4} (#2,#3);}
\draw[->] (0,-0.9) -- ++(0.5,0) -- ++(0,-0.35) node[right, align=left, xshift=1pt, yshift=6pt]{pick $N_i$; $PID=H_2(\lambda_V\|t)$;\\ $C_1=\mathrm{AEAD}_{k^\star}(PID\|N_i\|t)$; $t\leftarrow t+1$} -- ++(-0.5,0);
\msg{0}{4.2}{-2.0}{$M_1=\langle PID, C_1\rangle$}
\msg{4.2}{7.6}{-2.6}{resolve($PID$)}
\msg{7.6}{10.6}{-3.2}{window / replay check}
\ret{10.6}{7.6}{-3.8}{}
\ret{7.6}{4.2}{-4.3}{$(V,t)$}
\msg{4.2}{10.6}{-4.9}{is\_blocked($V$)?}
\draw[->] (4.2,-5.4) -- ++(0.5,0) -- ++(0,-0.35) node[right, align=left, xshift=1pt, yshift=4pt]{decrypt $C_1$ with $k^\star_V$; check $PID$, $t$} -- ++(-0.5,0);
\msg{4.2}{7.6}{-6.4}{accept($V,t,N_i$)}
\msg{7.6}{10.6}{-6.9}{accept\_epoch($V,t$)}
\draw[->] (4.2,-7.4) -- ++(0.5,0) -- ++(0,-0.35) node[right, align=left, xshift=1pt, yshift=4pt]{pick $N_j$; $K=\mathrm{KDF}(k^\star\|N_i\|N_j\|t)$} -- ++(-0.5,0);
\msg{4.2}{0}{-8.3}{$M_2=\langle ID_R,\mathrm{AEAD}_{k^\star}(N_j\|t;\,PID\|N_i)\rangle$}
\draw[->] (0,-8.8) -- ++(0.5,0) -- ++(0,-0.35) node[right, align=left, xshift=1pt, yshift=4pt]{verify $M_2$; derive the same $K$} -- ++(-0.5,0);
\end{tikzpicture}
\caption{Sequence diagram of one authentication}
\label{fig:seq}
\end{figure}

\subsection{Block Diagram}
Figure~\ref{fig:block} shows the blocks of the system and the flow between them: the offline enrolment, the online exchange between a device (Python or ESP8266) and the server with its pluggable resolver, and the evaluation pipeline that turns measured probe counts into the simulation results, figures and tables.

\begin{figure}[H]
\centering
\resizebox{\textwidth}{!}{%
\begin{tikzpicture}[font=\scriptsize, >=Stealth,
  blk/.style={draw, rounded corners=2pt, fill=#1, align=center, minimum height=10mm, text width=3.0cm},
  blk/.default=blue!5]
\node[blk=gray!10] (ta) at (0,0) {Trusted Authority\\ (\texttt{hea-ta enroll})};
\node[blk=gray!10] (cred) at (4.5,0) {Credential files\\ $\lambda$, peer secrets};
\node[blk, text width=3.2cm] (dev) at (0,-2.4) {Device\\ Python \texttt{hea-device}\\ or ESP8266 C core};
\node[blk, text width=3.6cm] (srv) at (8.2,-2.4) {Server \texttt{hea-server}\\ Resolver (O1--O4, TM) $\to$\\ AEAD verify $\to$ session log};
\draw[->] (ta) -- (cred);
\draw[->] (cred.south) -- ++(0,-0.5) -| (dev.north);
\draw[->] (cred.south) -- ++(0,-0.5) -| (srv.north);
\draw[->, thick] ([yshift=2.5mm]dev.east) -- node[above]{$M_1$ over TCP} ([yshift=2.5mm]srv.west);
\draw[<-, thick] ([yshift=-2.5mm]dev.east) -- node[below]{$M_2$} ([yshift=-2.5mm]srv.west);
\node[blk=green!8] (rc) at (0,-5.2) {\texttt{resolution\_costs.py}\\ \texttt{rfid\_schemes.py}};
\node[blk=green!8] (cm) at (4.1,-5.2) {\texttt{measured.csv}\\ \texttt{cost\_model.csv}};
\node[blk=green!8] (om) at (8.2,-5.2) {OMNeT++ ResolveNet\\ (960 runs, \texttt{results.csv})};
\node[blk=green!8] (mf) at (8.2,-7.4) {\texttt{make\_figs.py}\\ figures, ranked table};
\node[blk=green!8] (pv) at (0,-7.4) {ProVerif models\\ $\to$ logs};
\node[blk=gray!10] (rep) at (4.1,-7.4) {Paper and\\ this report};
\draw[->] (srv.south) -- ++(0,-0.6) -| node[pos=0.25, above]{resolvers and probe counters} (rc.north);
\draw[->] (rc) -- (cm); \draw[->] (cm) -- (om); \draw[->] (om) -- (mf);
\draw[->] (mf) -- (rep); \draw[->] (pv) -- (rep);
\end{tikzpicture}}
\caption{Block diagram of the system}
\label{fig:block}
\end{figure}

\subsection{Data Flow Diagram}
Figure~\ref{fig:dfd0} is the context-level (level~0) DFD: the system exchanges enrolment data with the trusted authority, requests and responses with devices, and receives forged or replayed messages from an adversary. Figure~\ref{fig:dfd1} refines it into level~1 processes and data stores.

\begin{figure}[H]
\centering
\begin{tikzpicture}[font=\footnotesize, >=Stealth,
  ext/.style={draw, fill=gray!10, minimum width=2.4cm, minimum height=9mm, align=center},
  proc/.style={draw, circle, fill=blue!6, align=center, minimum size=3.0cm}]
\node[proc] (p) at (0,0) {\textbf{0}\\ Pseudonym-based\\ authentication\\ system};
\node[ext] (ta) at (-5.2,1.8) {Trusted\\ Authority};
\node[ext] (dv) at (-5.2,-1.8) {Device};
\node[ext] (ad) at (5.2,-1.8) {Adversary};
\node[ext] (rs) at (5.2,1.8) {Researcher};
\draw[->] (ta) -- node[above, sloped]{credentials} (p);
\draw[->] ([yshift=1.5mm]dv.east) -- node[above, sloped]{$M_1$} ([yshift=1.5mm]p.215);
\draw[<-] ([yshift=-1.5mm]dv.east) -- node[below, sloped]{$M_2$, $K$} ([yshift=-1.5mm]p.215);
\draw[->] (ad) -- node[above, sloped]{forged / replayed $M_1$} (p);
\draw[<-] ([yshift=1.5mm]rs.west) -- node[above, sloped]{probe counts, latency} ([yshift=1.5mm]p.30);
\end{tikzpicture}
\caption{Data flow diagram, level 0 (context)}
\label{fig:dfd0}
\end{figure}

\begin{figure}[H]
\centering
\begin{tikzpicture}[font=\scriptsize, >=Stealth,
  ext/.style={draw, fill=gray!10, minimum width=1.9cm, minimum height=8mm, align=center},
  proc/.style={draw, circle, fill=blue!6, align=center, minimum size=1.9cm, inner sep=1pt},
  store/.style={draw=none, minimum width=2.8cm, minimum height=6mm, align=left, fill=white}]
\node[ext] (ta) at (-6.6,3.2) {Trusted\\Authority};
\node[proc] (p1) at (-3.0,3.2) {1.0\\ Enrol\\ entities};
\node[store] (d1) at (1.6,3.4) {D1\quad Credentials};
\node[ext] (dv) at (-6.6,0) {Device};
\node[proc] (p2) at (-3.0,0) {2.0\\ Build $M_1$};
\node[proc] (p3) at (1.2,0) {3.0\\ Resolve\\ pseudonym};
\node[store] (d3) at (3.6,1.8) {D3\quad Session log};
\node[proc] (p4) at (5.6,0) {4.0\\ Verify $M_1$,\\ build $M_2$};
\node[store] (d2) at (2.6,-2.6) {D2\quad Resolution table / state};
\node[ext] (ad) at (-0.8,-2.6) {Adversary};
\node[proc] (p5) at (-3.0,-3.6) {5.0\\ Verify $M_2$,\\ derive $K$};
\foreach \n in {d1,d2,d3}{\draw (\n.north west) -- (\n.north east) (\n.south west) -- (\n.south east) (\n.north west) -- (\n.south west);}
\draw[->] (ta) -- node[above]{IDs, master $s$} (p1);
\draw[->] (p1) -- node[above]{$\lambda$, peer secrets} (d1);
\draw[->] (dv) -- node[above]{request} (p2);
\draw[->] (d1.south west) -- node[left, pos=0.35]{$\lambda_V$, $k^\star$} (p2.north east);
\draw[->] (d1.south east) -- node[right, pos=0.3]{$k^\star$} (p4.north);
\draw[->] (p2) -- node[above]{$M_1$} (p3);
\draw[->] (ad) -- node[left, pos=0.4]{forged $PID$} (p3.south west);
\draw[<->] (p3.south east) -- node[right, pos=0.6]{lookup / probe} (d2.north);
\draw[->] (d3.south west) -- node[above left, pos=0.3]{window} (p3.north east);
\draw[->] (p3) -- node[below]{$(V,t)$, $M_1$} (p4);
\draw[->] (p4.north west) -- node[above right, pos=0.6]{accept, failures} (d3.south east);
\draw[->] (p4.south) -- ++(0,-4.2) -| node[pos=0.25, below]{$M_2$} (p5.south);
\draw[->] (p5.west) -| node[pos=0.75, left]{$K$} (dv.south);
\end{tikzpicture}
\caption{Data flow diagram, level 1}
\label{fig:dfd1}
\end{figure}

\subsection{Timeline Chart}
Figure~\ref{fig:timeline} lists the confirmed milestones supplied by the team. Unlike the activity bars in Figure~\ref{fig:gantt}, these dates are not estimates and do not depend on repository history. The events are spaced for readability, not in proportion to elapsed time.

\begin{figure}[H]
\centering
\begin{tikzpicture}[font=\small, >=Stealth]
\draw[very thick, ->] (0,0.6) -- (0,-8.2);
\fill[orange!80!black] (0,-0.0) circle (2.5pt);
\node[left, align=right] at (-0.2,-0.0) {\textbf{6 February}};
\node[right, text width=8.7cm, align=left] at (0.3,-0.0) {Base paper published};
\fill[orange!80!black] (0,-1.2) circle (2.5pt);
\node[left, align=right] at (-0.2,-1.2) {\textbf{5 March}};
\node[right, text width=8.7cm, align=left] at (0.3,-1.2) {Synopsis submitted};
\fill[orange!80!black] (0,-2.4) circle (2.5pt);
\node[left, align=right] at (-0.2,-2.4) {\textbf{17 March}};
\node[right, text width=8.7cm, align=left] at (0.3,-2.4) {Review 6.1};
\fill[orange!80!black] (0,-3.6) circle (2.5pt);
\node[left, align=right] at (-0.2,-3.6) {\textbf{15 April}};
\node[right, text width=8.7cm, align=left] at (0.3,-3.6) {Review 6.2};
\fill[orange!80!black] (0,-4.8) circle (2.5pt);
\node[left, align=right] at (-0.2,-4.8) {\textbf{28 April}};
\node[right, text width=8.7cm, align=left] at (0.3,-4.8) {Semester VI logbook and attendance report prepared};
\fill[orange!80!black] (0,-6.0) circle (2.5pt);
\node[left, align=right] at (-0.2,-6.0) {\textbf{21 August}};
\node[right, text width=8.7cm, align=left] at (0.3,-6.0) {Review 7.1};
\fill[orange!80!black] (0,-7.2) circle (2.5pt);
\node[left, align=right] at (-0.2,-7.2) {\textbf{27 October}};
\node[right, text width=8.7cm, align=left] at (0.3,-7.2) {Review 7.2};
\end{tikzpicture}
\caption{Confirmed project milestones in 2026 (spacing not to scale)}
\label{fig:timeline}
\end{figure}

\section{Hardware Requirements}
\begin{table}[H]
\centering
\caption{Hardware requirements}
\label{tab:hw}
\small
\renewcommand{\arraystretch}{1.2}
\begin{tabularx}{\textwidth}{|p{0.27\textwidth}|X|p{0.17\textwidth}|}
\hline
\textbf{Component} & \textbf{Specification / use} & \textbf{Status}\\ \hline
Development and evaluation laptop & Apple MacBook Pro (14-inch, 2023), identifier \texttt{Mac14,9}, model \texttt{MPHE3HN/A}. Apple M2 Pro: 10-core CPU (6 performance + 4 efficiency), 16-core GPU, 16\,GB unified memory, 512\,GB-class SSD. Model specifications include 200\,GB/s memory bandwidth and a 14.2-inch Liquid Retina XDR display (3024\,$\times$\,1964 pixels)~\cite{apple}. Used for the original OMNeT++ and ProVerif runs. & Already available\\ \hline
Machine used for this report's re-runs & 12th Gen Intel Core i5-12500H, 16 logical CPUs, 15\,GiB RAM, Linux 6.18 (tests, C vector check, interoperability, desynchronisation, probe sweep and figure generation) & Already available\\ \hline
ESP8266 development boards (NodeMCU), two units & Two team-owned boards were used to test the device-side C firmware: 80\,MHz stock clock (160\,MHz selectable), Wi-Fi, no hardware cryptography; epoch persisted to flash. The host build also checks the port without hardware. & Already owned; \rupee\,0 additional cost\\ \hline
Network & Wi-Fi access point / LAN connecting the board and the server laptop & Already available\\ \hline
\end{tabularx}
\end{table}

The team confirmed the configuration directly from the MacBook and confirmed that the same machine was used for the original OMNeT++ and ProVerif runs. Its current operating system is macOS~26.6.2 (build \texttt{25G83}). The 512\,GB-class SSD is reported by macOS as a 494.4\,GB container, with approximately 460\,GiB shown for the root filesystem; these are capacity-reporting values for the same drive. The Linux machine listed separately is the environment used for this report's re-runs.

\section{Software Requirements}
\begin{table}[H]
\centering
\caption{Software requirements}
\label{tab:sw}
\small
\renewcommand{\arraystretch}{1.2}
\begin{tabularx}{\textwidth}{|p{0.27\textwidth}|X|}
\hline
\textbf{Software} & \textbf{Version and purpose}\\ \hline
Operating system & MacBook current configuration: macOS~26.6.2 (build \texttt{25G83}); report re-runs: Linux~6.18; CI: Ubuntu\\ \hline
Python & 3.11 or newer (CI: 3.11 and 3.12; this report's re-runs: 3.14.9); reference implementation, measurement and figure scripts\\ \hline
\texttt{cryptography} & $\geq$\,41.0 (50.0.2 in the re-runs): AES-256-GCM, HKDF-SHA256, SHA-256, P-256\\ \hline
pytest, pytest-cov, Hypothesis & Unit, property and adversarial tests; coverage gate of 90\%\\ \hline
matplotlib & Figures (paper and this report)\\ \hline
C compiler (GCC/Clang), Make & Host build of the C device core and the vector check\\ \hline
BearSSL (ESP8266 Arduino core fork) & AES-256-GCM, SHA-256 and HKDF in the C device core\\ \hline
PlatformIO, Arduino framework & Building and flashing the ESP8266 firmware\\ \hline
OMNeT++ & 6.4.0: network simulation (ResolveNet, 960 runs)\\ \hline
ProVerif & 2.05 (built from source with OCaml 5.5.2): symbolic verification\\ \hline
Git, GitHub, GitHub Actions & Version control, code review through pull requests, continuous integration\\ \hline
\LaTeX{} (Tectonic 0.17.0, XeLaTeX) & Paper and this report\\ \hline
\end{tabularx}
\end{table}

\section{Cost Estimation}\label{sec:cost}
\begin{table}[H]
\centering
\caption{Cost estimation}
\label{tab:cost}
\small
\renewcommand{\arraystretch}{1.25}
\begin{tabularx}{\textwidth}{|>{\centering\arraybackslash}p{0.06\textwidth}|p{0.18\textwidth}|X|p{0.17\textwidth}|}
\hline
\textbf{Sr. No.} & \textbf{Item / Resource} & \textbf{Description} & \textbf{Estimated Cost (\rupee)}\\ \hline
1 & Hardware & Existing team laptops and two already-owned ESP8266 NodeMCU boards used for testing & \rupee\,0 additional cost\\ \hline
2 & Software / Tools & Python and libraries, GCC, BearSSL, PlatformIO, ProVerif, OMNeT++ (free for academic use), \LaTeX & \rupee\,0 (free/open-source)\\ \hline
3 & Repository / CI & GitHub repository and GitHub Actions (free tier) & \rupee\,0\\ \hline
4 & APIs / Services & None used & \rupee\,0\\ \hline
5 & Other Expenses & Publication or conference fees, printing and binding of this report & \textit{To be confirmed}\\ \hline
 & \textbf{Total Estimated Cost} & & \textit{To be confirmed}\\ \hline
\end{tabularx}
\end{table}
The estimate follows the department's instructions: resources that were already available (the team's laptops and the network) are listed as existing resources with no purchase cost, and all software is free or open source and is listed at \rupee\,0. Both ESP8266 boards were already owned by the team, so no board purchase was required. The confirmed additional cost of the development and evaluation resources is zero. Publication, printing and binding expenses, and therefore the overall estimated total, remain to be confirmed.

\section{Conclusion}
This chapter set out a methodology in which every phase builds on a checked result: the protocol is implemented from Section~III of~\cite{base} and pinned down with independent vectors; the need for resolution is derived from the specification alone; seven strategies drawn from RFID private identification are implemented with closed-form costs; their probe counts are measured, validated and priced under the paper's own constants in an OMNeT++ model of its network; and the security properties are checked in ProVerif. Requirements were gathered by analysing the specification and recorded in a conformance matrix and an assumption register. The layered architecture separates primitives, protocol logic, storage, transport and evaluation, so that resolvers can be exchanged without touching the protocol. All software is free or available for academic use, and the laptops and both development boards were already owned. The retrospective schedule records the confirmed milestones and approximate activity periods, providing context for the implementation and results in the next chapter.
